\section{Recording identity, not person ID alone, is a prerequisite for EEG transfer validation}
\label{sec:v26-identity}
Contemporary EEG reproducibility work already documents that
trial-level validation can greatly overestimate subject-independent
performance \cite{gil2026subjectvalidation}; neither the
principle of leakage-safe splitting nor signal hashing is
presented here as novel. The previously unverified
\emph{record identity across different IDs} and its
consequences for this frozen regularisation comparison
are the dataset-specific methodological findings.

\paragraph{An independent raw-data integrity check.}
A separate source-locked Liu2024 acute-stroke experiment
originally included 50 public subject identifiers and
2,000 nominal trials (40 per identifier), with
29 EEG channels and no target-statistic fitting for
its \emph{predeclared} source regularisation rule.
The frozen 50-identifier outcome means were 0.59400
AUC for source-macro selection and 0.58985 for
genuine source-trained-mean Gaussian selection.
These are not yet estimates of transfer to 50
independent processed recordings.

We subsequently inspected each public MOABB 1.7.2
processed EEG trial array and the associated binary
left/right event labels \emph{without examining classifier
scores as a criterion for record grouping}. Exact
float64 array hashes, verified over all 50 identifiers,
found that IDs 4 and 5 have identical entire processed
EEG trial arrays and identical class labels.
No other full processed EEG array was repeated,
so the fixed release provides 49 unique preprocessed
recording groups under this exact identity criterion.
Although IDs 47 and 49 had coincident vectors of
reported target AUCs, their EEG signals differed
and they must remain distinct.
An independent released-EDF inspection found that
IDs 4 and 5 were in different binary EDF files but
32 of their 33 raw numerical channels were identical;
the sole differing channel was index 32 (zero-based),
with an empty name in the MNE EDF reader.
The comparison was performed on numerical signals,
not merely identical file names or matched classifier scores.
These tests do not identify human patients and do
not establish why the release contains identical
EEG waveforms under two subject IDs.

\paragraph{Why this can invalidate apparent held-person isolation.}
Let $D_i$ denote the preprocessed trial-signal and
label array associated with nominal subject ID $i$.
In the original leave-one-ID-out training protocol,
$\mathcal S_i=\{D_k:k\ne i\}$. For $D_i=D_j$ with
$i\ne j$, $D_j\in\mathcal S_i$ is a copy of the
nominal held-out record, including its labels.
The source-fitted tangent reference, spectral covariance,
class contrast, Gaussian predictive source criterion and
discriminative source validation can therefore all
depend on the held record through its duplicate.
The assumption of independent source/target
\emph{recordings} is violated even though the numerical
subject identifiers differ. For all other outer folds,
the same recording receives additional weight through
its duplicate IDs. Correcting only the denominator
of a reported mean AUC does not repair the fitted
source models.

\paragraph{Purged re-fitting.}
The precommitted retrospective V26 protocol groups
only exact signal- and label-identical records;
it retains the smallest ID per group, excludes
other identical copies \emph{from every source pool
and every target denominator}, and re-fits all
source references and source-only selection rules.
Every model uses one of the other 48 unique
processed recording groups as its held target
and is trained on the remaining 48.
All grouped source decisions are locked before
their corresponding target AUCs are joined.
Numerical results are not reported until the
real-data run and identity receipts have passed.

\paragraph{Information-budget limitation.}
Identifying duplicates required a global inspection
of the publicly released held-target EEG arrays;
therefore the new reanalysis is \emph{not} a strict
no-target-feature-access Z0 or an independent
prospective test. This retrospective integrity check
cannot prove causal performance inflation solely
from the duplicate: purging also changes source
sample weights and fitted covariance/geometry.
The previously exposed 50-identifier result
must be treated as \emph{unqualified} with respect
to recording-independent transfer until a
record-level separation analysis is complete.
